\subsection{Derivations in separable extensions}

We have seen (V, p.~122, Th. 1) that every extension L of a field K of characteristic 0 is separable; moreover, every derivation of K into a vector space over L extends then to a derivation of L (V, p.~130, Cor. 1). More generally we have the following statement:

THEOREM 3. --- Let K be a field and L an extension of K. For L to be separable over K it is necessary and sufficient that every derivation of K into a vector L-space should extend to a derivation of L.

We may assume that K has characteristic \(\neq0\) . Suppose first that L is separable over K. Let V be a vector space over L and \(\Delta\) a derivation of K into V. By Mac

Lane's criterion (V, p.~124, Remark), the fields \(L^p\) and K are linearly disjoint over \(K^p\) . Since \(\Delta\) is a \(K^p\) -linear mapping of K into V, it extends in a unique way to an \(L^p\) -linear mapping \(\Delta'\) of \(K[L^p] = K(L^p)\) into V. It is clear that \(\Delta'\) is a derivation of \(K(L^p)\) into V which vanishes on \(L^p\) ; thus it extends (V, p.~101, Cor. 1) to a derivation of L into V.

Conversely, suppose that every derivation of K with values in L extends to a derivation of L into L. Let B be a p-basis of K (V, p.~99) and let \(\Lambda\) be the set of families \((\alpha_{b})_{b\in\mathbf{B}}\) with finite support consisting of integers between 0 and p-1. For each \(b\in B\) there exists a derivation \(\Delta_{b}\) of K into K characterized by \(\Delta_{b}(b')=\delta_{bb'}\) (Kronecker symbol) for each \(b'\in B\) (V, p.~103). By hypothesis there exists for each \(b\in B\) a derivation \(D_{b}\) of L into L extending \(\Delta_{b}\) . We have \(\mathrm{D}_{b}(b')=\delta_{bb'}\) for b, \(b'\) in B, which proves that in \(\Omega_{L^{p}}(L)\) the differentials db (for \(b\in B\) ) are linearly independent over L. Therefore (V, p.~102, Th. 1), B is p-free over \(L^{p}\) . It follows (V, p.~98, Prop. 1 and p.~124, Remark) that the extension L of K is separable.

COROLLARY. --- If L is a separable extension of K, the canonical mapping \(\alpha_{\mathrm{P}}:\Omega_{\mathrm{P}}(\mathrm{K})\otimes_{\mathrm{K}}\mathrm{L}\to\Omega_{\mathrm{P}}(\mathrm{L})\) is injective for every subfield P of K. Conversely, if there exists a perfect subfield P of K (for example the prime subfield of K) such that the mapping \(\alpha_{P}\) is injective, then L is separable over K.

The first assertion follows from Prop. 2 (V, p.~128) and Th. 3. Conversely, let P be a perfect subfield of K; we have \(P = P^{p} \subset K^{p}\) , hence every derivation of K into a vector L-space is a P-derivation; the second assertion of the corollary now follows from Cor. 2 (V, p.~130) and Th. 3.

